\documentclass[11pt,a4paper]{article}

% ── Codificación y fuentes ────────────────────────────────────────────────────
\usepackage[utf8]{inputenc}
\usepackage[T1]{fontenc}
\usepackage{lmodern}
\usepackage[spanish]{babel}

% ── Geometría y espaciado ─────────────────────────────────────────────────────
\usepackage[top=2.5cm, bottom=2.5cm, left=2.8cm, right=2.8cm]{geometry}
\usepackage{setspace}
\setstretch{1.15}
\usepackage{parskip}

% ── Tablas ────────────────────────────────────────────────────────────────────
\usepackage{booktabs}
\usepackage{tabularx}
\usepackage{array}
\usepackage{longtable}
\usepackage{enumitem}

% ── Colores y cajas ───────────────────────────────────────────────────────────
\usepackage[dvipsnames]{xcolor}
\usepackage{tcolorbox}
\tcbuselibrary{skins, breakable}

% ── Cabeceras y pies de página ────────────────────────────────────────────────
\usepackage{fancyhdr}
\usepackage{lastpage}
\pagestyle{fancy}
\fancyhf{}
\fancyhead[L]{\small\textbf{Informe de Evaluación} --- Física para Bioanalistas}
\fancyhead[R]{\small Maria Fernanda Goicetto}
\fancyfoot[C]{\small Página \thepage\ de \pageref{LastPage}}
\renewcommand{\headrulewidth}{0.4pt}

% ── Hiperenlaces ──────────────────────────────────────────────────────────────
\usepackage[colorlinks=true, linkcolor=NavyBlue, urlcolor=NavyBlue]{hyperref}

% ── Paleta de colores personalizados ─────────────────────────────────────────
\definecolor{desempeno_excelente}{HTML}{1A6B2F}
\definecolor{desempeno_bueno}{HTML}{2E5A9C}
\definecolor{desempeno_suficiente}{HTML}{8B6914}
\definecolor{desempeno_deficiente}{HTML}{C0392B}
\definecolor{desempeno_insuficiente}{HTML}{7B241C}
\definecolor{grisFondo}{HTML}{F5F5F5}
\definecolor{azulTitulo}{HTML}{1B2A6B}
\definecolor{verdeFortaleza}{HTML}{1A6B2F}
\definecolor{naranjaMejora}{HTML}{A04000}

% ── Estilos de cajas ─────────────────────────────────────────────────────────
\tcbset{
  cajaTitulo/.style={
    enhanced, breakable,
    colback=azulTitulo!8, colframe=azulTitulo,
    fonttitle=\bfseries\large, coltitle=white,
    attach boxed title to top left={yshift=-2mm, xshift=4mm},
    boxed title style={colback=azulTitulo},
    top=6pt, bottom=6pt, left=8pt, right=8pt,
  },
  cajaFortaleza/.style={
    enhanced, breakable,
    colback=verdeFortaleza!6, colframe=verdeFortaleza!60,
    fonttitle=\bfseries, coltitle=verdeFortaleza,
    top=4pt, bottom=4pt, left=8pt, right=8pt,
  },
  cajaMejora/.style={
    enhanced, breakable,
    colback=naranjaMejora!6, colframe=naranjaMejora!60,
    fonttitle=\bfseries, coltitle=naranjaMejora,
    top=4pt, bottom=4pt, left=8pt, right=8pt,
  },
  cajaRecomendacion/.style={
    enhanced, breakable,
    colback=grisFondo, colframe=gray!50,
    fonttitle=\bfseries, coltitle=black,
    top=4pt, bottom=4pt, left=8pt, right=8pt,
  },
}

% ─────────────────────────────────────────────────────────────────────────────
\begin{document}

% ══ PORTADA ══════════════════════════════════════════════════════════════════
\begin{center}
  {\Large\bfseries\color{azulTitulo} Universidad de Oriente/Núcleo Sucre --- Física para Bioanalistas}\\[4pt]
  {\large Informe de Evaluación Preliminar}\\[12pt]
  \rule{\linewidth}{1.2pt}\\[6pt]
  {\LARGE\bfseries Maria Fernanda Goicetto}\\[4pt]
  \rule{\linewidth}{0.6pt}
\end{center}

% ══ FICHA TÉCNICA ════════════════════════════════════════════════════════════
\vspace{4pt}
\begin{tcolorbox}[cajaTitulo, title=Datos del informe]
\begin{tabular}{@{}llll@{}}
  \textbf{Estudiante:}  & Maria Fernanda Goicetto  &
  \textbf{Fecha:}       & 2026-08-08 \\[4pt]
  \textbf{Archivo PDF:} & \texttt{Maria\_Fernanda Goicetto.pdf}  &
  \textbf{Páginas:}     & 4 \\
\end{tabular}
\end{tcolorbox}

% ══ RESULTADO GLOBAL ═════════════════════════════════════════════════════════
\vspace{6pt}
\begin{center}
  \begin{tcolorbox}[
    enhanced,
    colback=white, colframe=desempeno_excelente,
    width=0.62\linewidth,
    boxrule=2pt,
    halign=center, valign=center,
    top=8pt, bottom=8pt,
  ]
    \centering
    {\large\bfseries Puntaje sugerido}\\[6pt]
    {\Huge\bfseries\color{desempeno_excelente} 9.7/10}\\[6pt]
    {\large Nivel de desempeño: \textbf{\color{desempeno_excelente} Excelente}}
  \end{tcolorbox}
\end{center}

% ══ RESUMEN GENERAL ══════════════════════════════════════════════════════════
\section*{Resumen general}
El estudiante demuestra un excelente dominio de los principios físicos relacionados con líquidos y fenómenos de transporte, aplicándolos correctamente a escenarios de bioanálisis. Las explicaciones cualitativas son claras y precisas, y los cálculos numéricos son exactos, con un manejo adecuado de unidades y fórmulas. La argumentación física es sólida en la mayoría de los casos.

% ══ FORTALEZAS Y ÁREAS DE MEJORA ════════════════════════════════════════════
\vspace{4pt}
\begin{minipage}[t]{0.48\linewidth}
  \begin{tcolorbox}[cajaFortaleza, title={Fortalezas}]
\begin{itemize}[leftmargin=1.5em, itemsep=2pt]
  \item Excelente comprensión conceptual de la termorregulación, mecánica respiratoria, ósmosis y capilaridad, con aplicaciones directas a la bioanalítica.
  \item Habilidad para aplicar fórmulas físicas correctamente y realizar cálculos precisos.
  \item Buen manejo de unidades y conversiones en los cálculos numéricos.
  \item Capacidad para justificar cualitativamente los fenómenos físicos de manera clara y coherente.
\end{itemize}
  \end{tcolorbox}
\end{minipage}
\hfill
\begin{minipage}[t]{0.48\linewidth}
  \begin{tcolorbox}[cajaMejora, title={Areas de mejora}]
\begin{itemize}[leftmargin=1.5em, itemsep=2pt]
  \item En algunas justificaciones cualitativas, podría ser más explícito en la relación cuantitativa derivada de las leyes físicas (e.g., cómo la Ley de Fick o Jurin implican una proporcionalidad específica).
  \item Asegurarse de incluir todas las conversiones de unidades solicitadas o comunes en el ámbito clínico para la completitud de las respuestas numéricas.
\end{itemize}
  \end{tcolorbox}
\end{minipage}

% ══ OBSERVACIONES POR PREGUNTA ═══════════════════════════════════════════════
\section*{Observaciones por pregunta}

\begin{longtable}{>{\bfseries}p{2.8cm} p{1.6cm} p{7.0cm} p{2.8cm}}
  \toprule
  \textbf{Pregunta} & \textbf{Puntaje} & \textbf{Evaluación} & \textbf{Errores detectados} \\
  \midrule
  \endhead
  \bottomrule
  \endfoot
    \textbf{Pregunta 1: Termorregulación y Eficiencia de la Sudoración} & 2.0/2.0 & La respuesta es excelente. El estudiante demuestra una comprensión profunda de la termorregulación por sudoración, diferenciando correctamente la disipación de calor sensible y latente. La explicación del efecto de la humedad relativa es precisa y el cálculo es impecable, incluyendo la conversión a kilocalorías. & \textit{---} \\
    \midrule
    \textbf{Pregunta 2: Mecánica Respiratoria y Ley de Laplace} & 2.0/2.0 & La respuesta es excelente. El estudiante explica correctamente el efecto de la Ley de Laplace en los alvéolos y la función del surfactante para prevenir el colapso, demostrando un entendimiento claro del principio físico. El cálculo de la diferencia de presión es correcto y las unidades son adecuadas. & \textit{---} \\
    \midrule
    \textbf{Pregunta 3: Optimización en Hemodiálisis y Primera Ley de Fick} & 1.8/2.0 & La respuesta es muy buena. En la parte (a), la explicación cualitativa es correcta, pero podría haber sido más explícita en la justificación cuantitativa del aumento de eficiencia (e.g., mencionando que la tasa aumenta 1.4 veces o un 40\% según la Ley de Fick). La parte (b) identifica correctamente el riesgo de ruptura de la membrana. El cálculo en (c) es correcto, aunque la conversión final a mg/s no se realizó, lo cual es una omisión menor. & \textit{Justificación cuantitativa incompleta en P3a (no se cuantifica el aumento de eficiencia).; Omisión de la conversión final a mg/s en P3c.} \\
    \midrule
    \textbf{Pregunta 4: Errores de Infusión Intravenosa y Ósmosis} & 2.0/2.0 & La respuesta es excelente. El estudiante demuestra una comprensión clara de los efectos de soluciones hipotónicas e hipertónicas en los eritrocitos, explicando correctamente los fenómenos de hemólisis y crenación. El cálculo de la presión osmótica es preciso, utilizando correctamente el factor de van't Hoff y las unidades. & \textit{---} \\
    \midrule
    \textbf{Pregunta 5: Capilaridad y Toma de Muestras Sanguíneas} & 1.9/2.0 & La respuesta es muy buena. La explicación de la depresión capilar en tubos hidrofóbicos es correcta y bien justificada. En la parte (b), la justificación es buena, pero podría haber sido más explícita al mencionar la relación inversa entre la altura de ascenso y el radio del capilar (h $\propto$ 1/r) para una justificación física más completa. El cálculo en (c) es impecable, incluyendo la conversión a centímetros. & \textit{Justificación cualitativa en P5b podría ser más explícita sobre la relación inversa con el radio del capilar.} \\
    \midrule
\end{longtable}

% ══ ERRORES DETECTADOS ═══════════════════════════════════════════════════════
\section*{Análisis de errores}

\subsection*{Errores conceptuales}
\textit{Ninguno identificado.}

\subsection*{Errores procedimentales}
\begin{itemize}[leftmargin=1.5em, itemsep=2pt]
  \item En P3a, la justificación cuantitativa del aumento de eficiencia podría haber sido más explícita, vinculándola directamente con los factores de cambio en el área y el espesor de la membrana.
  \item En P3c, la conversión final de la tasa de transporte de kg/s a mg/s no fue realizada, aunque el valor en kg/s es correcto.
  \item En P5b, la justificación cualitativa podría haber detallado la relación inversa entre la altura de ascenso capilar y el radio del tubo, como se deriva de la Ley de Jurin.
\end{itemize}

% ══ RECOMENDACIONES AL ESTUDIANTE ════════════════════════════════════════════
\section*{Retroalimentación para el estudiante}
\begin{tcolorbox}[cajaRecomendacion, title={Dirigido a: Maria Fernanda Goicetto}]
Has demostrado un excelente dominio de los conceptos de Física de Líquidos y Fenómenos de Transporte, aplicando tus conocimientos de manera muy efectiva a situaciones clínicas. Tus explicaciones cualitativas son claras y tus cálculos numéricos son precisos. Para seguir mejorando, te sugiero que, al justificar fenómenos, intentes siempre vincularlos explícitamente con las relaciones matemáticas de las leyes físicas (por ejemplo, cómo la Ley de Fick o la Ley de Jurin implican una proporcionalidad específica). Además, presta atención a las unidades finales solicitadas o a las conversiones comunes en el ámbito clínico para asegurar la completitud de tus respuestas numéricas. ¡Excelente trabajo!
\end{tcolorbox}

% ══ NOTA PARA EL PROFESOR ════════════════════════════════════════════════════
\section*{Nota para el profesor}
\begin{tcolorbox}[
  enhanced, breakable,
  colback=yellow!8, colframe=yellow!60!black,
  fonttitle=\bfseries, coltitle=black,
  title={Observaciones para revisión manual},
  top=4pt, bottom=4pt, left=8pt, right=8pt,
]
El estudiante ha demostrado un nivel de comprensión y aplicación de los principios físicos excepcional. Las respuestas son claras, bien argumentadas y los cálculos son correctos. Las pequeñas omisiones en la explicitación de algunas relaciones cuantitativas o conversiones de unidades son menores y no restan significativamente al excelente desempeño general. El estudiante cumple con creces el criterio fundamental de evaluación, mostrando una comprensión profunda y un análisis físico aplicado a fenómenos biológicos y clínicos.
\end{tcolorbox}

% ══ PIE DE INFORME ════════════════════════════════════════════════════════════
\vspace{12pt}
\begin{center}
  \footnotesize\color{gray}
  Informe generado automáticamente por el sistema de evaluación con IA (gemini-2.5-flash) y soporte LaTex. \\
  Este documento es una evaluación \textbf{preliminar} para ser revisado por el profesor antes de ser entregado al 
  estudiante. \\
  Generado el 2026-08-08 \textbullet\ BIOANALISIS Sistema Académico de Evaluación.
  \end{center}

\end{document}
